In this section, we study proportionality under uniform random sampling (i.e., draw $|N|$ individuals i.i.d. from the uniform distribution on $\N$). In particular, we show that proportionality is well preserved under random sampling. This allows us to design efficient implementations of Algorithm~\ref{alg:ball} and Algorithm~\ref{alg:local}, and to introduce an efficient algorithm for auditing proportionality. %and introduce an efficient local search heuristic for searching for a proportional solution (Algorithm~\ref{alg:local}). 
We first present the general property and then demonstrate its various applications. 

\subsection{Proportionality Under Random Sampling}
For any $X \subseteq \M$ of size $k$ and center $y \in \M$, define $R(\N, X, y) = \{i \in \N : D_i(X) > \rho \cdot d(i,y) \}$. Note that solution $X$ is not $\rho$-proportional with respect to $\N$ if and only if there is some $y \in \M$ such that $\frac{|R(\N,X,y)|}{|\N|} \geq \frac{1}{k}$. 
%The error introduced by sampling is parameterized by $\epsilon$ and can be driven arbitrarily small by taking a larger sample. The error appears in the number of points necessary to constitute a blocking coalition. In particular, 
A random sample approximately preserves this fraction for all solutions $X$ and deviating centers $y$. The important idea in the proof is that we take a union bound over all possible solutions and deviations, and there are only $k {m \choose k}$ such combinations. 

\begin{theorem}
\label{thm:sampling}
Given $\N$, $\M$ and parameter $\rho \ge 1$, fix parameters $\epsilon, \delta \in [0,1]$. Let $N \subseteq \N$ of size $\Omega\left(\frac{k^3}{\epsilon^2} \log \frac{m}{\delta}\right)$ be chosen uniformly at random. Then, with probability at least $1-\delta$, the following holds for all $(X,y)$:
$$ \left| \frac{| R(N,X,y) |}{|N|} - \frac{| R(\N,X,y)  |}{|\N|} \right| \le \frac{\epsilon}{k}$$	
%	Let $N \subset \N$ be a uniform random sample of size $|N| = \frac{ 2 (1+\epsilon) k}{\epsilon^2} \ln \left(\frac{m}{1-\delta}\right)$. Suppose $X \subseteq M$ with $|X| = k$ is $\rho$-proportional with respect to $N$. Then with probability at least $1-\delta$, $X$ is $\rho$-proportional to $(1+\epsilon)$-deviations with respect to $\N$.
\end{theorem}
\begin{proof}
	Recall that $N$ is a random sample of $\N$. Hoeffding's inequality implies that for any fixed $(X,y)$, a sample of size $|N| = O\left( \frac{1}{\hat{\epsilon}^2} \log \frac{1}{\hat{\delta}} \right)$ is sufficient to achieve
$$  \left| \frac{| R(N, X,y) |}{|N|} - \frac{| R(\N, X, y)  |}{|\N|} \right| \le \hat{\epsilon}$$ 
with probability at least $1-\hat{\delta}$. Note that there are $q = m {m \choose k}$ possible choices of $(X,y)$ over which we take the union bound. Setting $\delta = \frac{\hat{\delta}}{q}$, and $\epsilon = \frac{\hat{\epsilon}}{k}$ is sufficient for the union bound to yield the theorem statement.
\end{proof}
\iffalse
\begin{proof}
The crucial idea of the proof is to bound the probability that we ``miss'' a blocking coalition by taking the union bound over possible centers, \textit{not} over exponentially many subsets of points. In particular, for each $y \in \M$, let $S_y$ be the $(1+\epsilon)\frac{n}{k}$ points $i \in \N$ with largest values for the ratio $\frac{D_i(X)}{d(i, y)}$. We will show that with high probability, $| N \cap S_y | \geq \frac{|N|}{k}$ for all $y \in \N$. Let 
$$\mu = \mathbb{E}\left[| N \cap S_y |\right] = \frac{|N|}{n} |S_y| = 2 \left(\frac{1+\epsilon}{\epsilon}\right)^2 \ln \left(\frac{m}{1-\delta}\right). $$ 
A standard application of the multiplicative Chernoff Bound yields 
$$\Pr \left( | N \cap S_y | < (1-\zeta) \mu \right) \leq e^{-\frac{\zeta^2 \mu }{2}}.$$  
Setting $\zeta = \frac{\epsilon}{1+\epsilon}$, 
$$\Pr \left( | N \cap S_y | < \frac{|N|}{k} \right) \leq e^{-\ln \left(\frac{m}{1-\delta}\right)} = \frac{1-\delta}{m}.$$
Taking the union bound over all $y \in \M$, we find that
$$\Pr \left( \exists y \in \M \mbox{ s.t. } | N \cap S_y | < \frac{|N|}{k} \right) \leq 1 - \delta. $$
Now, suppose that $X$ is $\rho$-proportional with respect to $N$, but is not $\rho$-proportional to $(1+\epsilon)$-deviations with respect to $\N$. Then there is some $y \in \M$ and $S \subseteq N$ of size $|S| \geq (1+\epsilon)\frac{n}{k}$ such that for all $i \in S$, $\rho \cdot d(i, y) < D_i(X)$. Then for \textit{any} subset $S' \subseteq S$, it remains true that for all $i \in S'$, $\rho \cdot d(i, y) < D_i(X)$. In particular, this is true for $S' = N \cap S_y$. Since $X$ is $\rho$-proportional with respect to $N$, it must be the case that $|N \cap S_y| < \frac{|N|}{k}$. But there exists such a $y$ with probability at most $1-\delta$, as we have shown.  
\end{proof}
\fi

In order to apply the above theorem, we say that a solution $X$ is $\rho$-proportional to $(1+\epsilon)$-deviations if for all $y \in \M$ and for all $S \subseteq \N$ where $|S| \geq (1+\epsilon) \frac{n}{k}$, there exists some $i \in S$ such that $\rho \cdot d(i, y) \geq D_i(X)$. Note that if $X$ is $\rho$-proportional to 1-deviations, it is simply $\rho$-proportional. We immediately have the following: 

\begin{corollary}
Let $N \subset \N$ be a uniform random sample of size $|N| = \Omega \left(\frac{k^3}{\epsilon^2} \ln \frac{m}{\delta}\right)$. Suppose $X \subseteq M$ with $|X| = k$ is $\rho$-proportional with respect to $N$. Then with probability at least $1-\delta$, $X$ is $\rho$-proportional to $(1+\epsilon)$-deviations with respect to $\N$.
\end{corollary}
\iffalse
\begin{proof}
Since $X$ is $\rho$-proportional with respect to $N$, we have $\frac{| Q(X,y) |}{|N|} \le \frac{1}{k}$ for all $y \in \M$. By Theorem~\ref{thm:sampling}, we have $\frac{| R(X,y)  |}{|\N|} \le \frac{1+\epsilon}{k}$ for all $y \in \M$ with high probability. But this means $X$ is $\rho$-proportional to $(1+\epsilon)$-deviations with respect to $\N$.
\end{proof}
\fi

\subsection{Linear Time Implementation}
We now consider how to take advantage of Theorem~\ref{thm:sampling} to optimize Algorithm~\ref{alg:ball} and Algorithm~\ref{alg:local}. First, note that Algorithm~\ref{alg:ball} takes $\tilde{O}(m n)$ time, which is quadratic in input size. % with a standard implementation using priority queues.  This becomes unsatisfactory when $m$ and $n$ are both very large. 
A corollary of Theorem~\ref{thm:sampling} is that we can approximately implement Algorithm~\ref{alg:ball} in nearly linear time, comparable to the running time of the standard $k$-means heuristic.

\begin{corollary} 
\label{cor:fastGreedy}
Algorithm~\ref{alg:ball}, when run on $\M$ and a random sample $N \subseteq \N$ of size $|N| = \tilde{\Theta}\left(\frac{k^3}{\epsilon^2}\right)$, provides a solution that is $(1+\sqrt{2})$-proportional to $(1+\epsilon)$-deviations with high probability in $\tilde{O}\left( \frac{k^3}{\epsilon^2} m \right)$ time.
\end{corollary}

We also get a substantial speedup for our Local Capture algorithm. Recall that Local Capture (Algorithm~\ref{alg:local}) is an iterative algorithm that takes a target value of $\rho$ as a parameter, and if it converges, returns a $\rho$-proportional clustering. Without sampling, each iteration of Algorithm~\ref{alg:local} takes $\tilde{O}(m n^2)$ time. Another corollary of Theorem~\ref{thm:sampling} is that it is sufficient to run the Local Capture on a random sample of $k^3/\epsilon^2$ out of the $n$ points in $\N$ in order to search for a clustering that is $\rho$-proportional with respect to $(1+\epsilon)$-deviations.
%Corollary~\ref{cor:fastGreedy} follows from running Algorithm~\ref{alg:ball} on $\M$ and a random sample $N \subseteq \N$ of size $|N| = \tilde{O}\left(\frac{k^3}{\epsilon^2}\right)$. Note that this speedup can be substantial. Consider the common special case where $\M = \N$ (i.e., there is only one set of data points given as input) and $n$ is large. For such problems, the $k$-means heuristic is frequently used for reasons of efficiency, as each iteration takes $O(kn)$ time. For constant $\epsilon$, our entire optimized Algorithm~\ref{alg:ball} runs in $\tilde{O}(k^3 n)$ time, so in essence, our algorithm can scale to any data set for which the $k$-means heuristic scales well.

\subsection{Efficient Auditing}
Alternatively, one might still want to run a non-proportional clustering algorithm, and ask whether the solution produced happens to be proportional. We call this the \textit{Audit Problem}. Given $\N, \M,$ and $X \subset \N$ with $|X| \leq k$, find the minimum value of $\rho$ such that $X$ is $\rho$-proportional. %If $\rho > 1$, find a set $S \subseteq \N$ of size $|S| \geq \frac{n}{k}$ and a point $y \in \M$ such that for all $i \in S, \rho \cdot d(i, y) < D_i(X)$. 
It is not too hard to see that one can solve the Audit Problem exactly in $O((k+m)n)$ time by computing for each $y \in \M$, the quantity $\rho_y$, the $\lceil \frac{n}{k} \rceil$ largest value of $\frac{D_i(X)}{d(i, y)}$. We subsequently find the $y$ that maximizes $\rho_y$. Again, this takes quadratic time, which can be worse than the time taken to find the clustering itself.

%Again, the dependence on the product of $m$ and $n$ (or $n^2$ in the common special case where $\M = \N$) is unsatisfactory. In particular, one would like an audit algorithm that takes no longer than the actual clustering algorithm, and as we have already mentioned, popular heuristics like the $k$-means heuristic run in $O(k n)$ time (for a constant number of iterations).


\iffalse
\begin{algorithm}
\begin{algorithmic}[1]
\caption{Audit}
\label{alg:audit}
	\FOR{$y \in \M$}
		\STATE Find $\rho_y$, the $\lceil \frac{n}{k} \rceil$ largest value of $\frac{D_i(X)}{d(i, y)}$, using linear time selection.
	\ENDFOR 
	\STATE $y^* \gets \mbox{argmax}_y \rho_y$
	\STATE Find $S$, the $\lceil (1-\epsilon) \frac{n}{k} \rceil$ points $i \in \N$ with largest $\frac{D_i(X)}{d(i, y^*)}$.
	\STATE Return $(\rho_{y^*}, S)$.
\end{algorithmic}
\end{algorithm}

Note that the reconstruction of $S$ can be done in $O(n)$ time, because we do not need to return $S$ in any particular order. Again, the dependence on the product of $m$ and $n$ (or $n^2$ in the common special case where $\M = \N$) is unsatisfactory. In particular, one would like an audit algorithm that takes no longer than the actual clustering algorithm, and as we have already mentioned, popular heuristics like the $k$-means heuristic run in $O(k n)$ time (for a constant number of iterations).
\fi

%It is not too hard to see that one can solve the Audit Problem exactly as follows. Scan over every possible center $y \in \M$ and consider the ratio $\frac{D_i(X)}{d(i, y)}$ for $i \in \N$. For every $y \in \M$, find the $\frac{n}{k}$ largest such ratio. The final solution $\rho$ is the maximum of such values found over all $y \in \M$ and 1. If the max is greater than 1, then the $\lceil \frac{n}{k} \rceil$ points with largest ratios for that $y$ constitute the blocking coalition. Implemented via linear time selection, the naive audit algorithm runs in $O((k+m)n)$ time. Again, the dependence on the product of $m$ and $n$ (or $n^2$ in the common special case where $\M = \N$) is unsatisfactory. In particular, one would like an audit algorithm that takes no longer than the actual clustering algorithm, and as we have already mentioned, popular heuristics like the $k$-means heuristic run in $O(k n)$ time (for a constant number of iterations). 

Consider a slightly relaxed $(\epsilon, \delta)$-Audit Problem where we are asked to find the minimum value of $\rho$ such that $X$ is $\rho$-proportional to $(1+\epsilon)$-deviations with probability at least $1-\delta$. %and if $\rho > 1$, we should find some $S \subseteq \N$ of size at least $|S| \geq (1-\epsilon) \frac{n}{k}$ and a point $y \in \M$ such that for all $i \in S, \rho \cdot d(i, y) < D_i(X)$. The algorithm is allowed to fail (report an an incorrect $\rho$ or $S$) with probability $\delta$. 
This problem can be efficiently solved by using a random sample $N \subseteq \N$ of points to conduct the audit.

\begin{corollary}
\label{cor:fastAudit}
The $(\epsilon, \delta)$-Audit Problem can be solved in $\tilde{O}\left( \left(k + m \right)\frac{k^3}{\epsilon^2} \right)$ time.
\end{corollary}  

%Corollary~\ref{cor:fastAudit} follows from running Algorithm~\ref{alg:audit} on $\M$ and a random sample $N \subseteq \N$ of size $|N| = \frac{ 2 (1+\epsilon) k}{\epsilon^2} \ln \left(\frac{m}{1-\delta}\right)$, except that one still reconstructs $S$ as the $(1-\epsilon) \frac{n}{k}$ points with largest ratio from the original set $|\N$ using linear time selection. The correctness follows easily from Theorem~\ref{thm:sampling}. 


%\begin{proof}
%	For each possible center $y \in \N$ calculate $\frac{D_i(X)}{d(i, y)}$ for every $i \in \N$. Sort these values in non-increasing order. Record the $\lceil \frac{n}{k} \rceil$ largest such ratio in this order. The final solution $\rho$ is the maximum of such values found over all $y \in \N$ and 1. If the max is greater than 1, then the $\lceil \frac{n}{k} \rceil$ points with largest ratios for that $y$ constitute the blocking coalition. 
	
%	It is clear that this solution does not underestimate the minimum value of $\rho$ such that $X$ is a $\rho$-core, since it only reports a $\rho > 1$ if it provides an accompanying coalition such that decreasing $\rho$ any further would violate the $\rho$-core condition. Suppose instead for a contradiction that this solution overestimates $\rho$. Then there exists some $y \in \N$ and $S^* \subseteq \N$ of size at least $\lceil \frac{n}{k} \rceil$ such that $\rho \cdot d(i, y) < D_i(X)$ for all $i \in S$, where $\rho$ is the value returned by the algorithm. Then the $\lceil \frac{n}{k} \rceil$ largest value of $\frac{D_i(X)}{d(i, y)}$ is strictly greater than $\rho$. But then our algorithm would have returned a value strictly greater than $\rho$, a contradiction. 
%\end{proof}

%For small clustering instances, this exact solution to the Audit Problem is sufficient. However, it does require access to all pairwise distances. Many popular clustering algorithms, and most notably the $k$-means algorithm, are very efficient for large data sets precisely because they do not require the computation of the entire matrix of pairwise distances. In such instances, solving the audit problem exactly may take longer than solving the original clustering problem, which is somewhat unsatisfactory. We therefore ask whether it is possible to approximately solve the Audit Problem in time that is asymptotically dominated by $mn$. In particular, we give a sampling based algorithm that has only logarithmic dependence on $n$, the size of the data set to be audited. First, we need to define the notion of an approximate audit.

% BEGIN COMMENT 
\iffalse

\begin{definition}[$\epsilon$-Approximate Audit Problem]
\label{def:approx-audit}
	Given $\N, \M$, and $X \subset \N$ with $|X| \leq k$, find a value of $\rho$ such that for all $y \in \M$ and for all $S \subseteq \N$ where $|S| \geq (1+\epsilon) \lceil \frac{n}{k} \rceil$, there exists some $i \in S$ such that $\rho \cdot d(i, y) \geq D_i(X)$. If $\rho > 1$, find a coalition $S \subseteq \N$ of size $|S| \geq (1-\epsilon) \lceil \frac{n}{k} \rceil$ and a point $y \in \M$ such that for all $i \in S, \rho \cdot d(i, y) < D_i(X)$.
\end{definition} 

Intuitively, the $\epsilon$-Approximate Audit Problem is the same as the Audit Problem except that we only guarantee to find proportionality violations of slightly more than $\lceil \frac{n}{k} \rceil$ points, and we only require the audit to find slightly fewer than $\lceil \frac{n}{k} \rceil$ points to support a claim of such a violation. We show that approximate auditing is amenable to random sampling. In particular, the sampling procedure is much simpler than the celebrated $k$-means++ sampling method for initializing the $k$-means algorithm~\cite{kmeans++}. Intuitively, this is because only large groups of cohesive points are entitled under proportionality, and outliers cannot therefore unduly affect the preservation of proportionality under sampling.

\begin{algorithm}
\begin{algorithmic}[1]
\caption{Audit}
\label{alg:audit}
	\FOR{$y \in \M$}
		\STATE Find $\rho_y$, the $\lceil \frac{\log(n)}{\epsilon^2} \rceil$ largest value of $\frac{D_i(X)}{d(i, y)}$, using linear time selection.
	\ENDFOR 
	\STATE $y^* \gets \mbox{argmax}_y \rho_y$
	\STATE Find $S$, the $\lceil (1-\epsilon) \frac{n}{k} \rceil$ points $i \in \N$ with largest $\frac{D_i(X)}{d(i, y^*)}$.
	\STATE Return $(\rho_{y^*}, S)$.
\end{algorithmic}
\end{algorithm}

\begin{theorem}
\label{thm:audit}
	Algorithm~\ref{alg:audit} solves the $\epsilon$-Approximate Audit Problem with high probability in $$O \left( \left( \frac{k \, m}{\epsilon^{2}} + n \right) \log(n) \right) \mbox{ time.}$$
\end{theorem}   
\begin{proof}
	\textbf{[Sketch of proof needs to be filled out]}. The run time is dominated by sorting the sample of $\frac{k \, log(n)}{\epsilon^2}$ points $n$ times, once for each possible $y \in \M$. The correctness of Algorithm~\ref{def:approx-audit} follows from concentration bounds. In particular, for each $y \in \M$, let $S_y$ be the $\lceil (1+\epsilon) \frac{n}{k} \rceil$ points $i \in \N$ with largest values for the ratio $\frac{D_i(X)}{d(i, y)}$. We will show that with high probability, $| M \cap S_y | \geq \lceil \frac{\log(n)}{\epsilon^2} \rceil$ for all $y \in \N$, where $M$ is the sample from Algorithm~\ref{alg:audit}. This follows from a Chernoff bound and union bound, and is sufficient to prove the value of $\rho$ returned by Algorithm~\ref{alg:audit} is not an underestimate (in the sense of Definition~\ref{def:approx-audit}). 
	
	A separate Chernoff bound shows that for each $y \in \N$, there are at least $(1-\epsilon) \lceil \frac{n}{k} \rceil$ points $i \in \N$ for which $\frac{D_i(X)}{d(i, y)}$ is at least the $\lceil \frac{\log(n)}{\epsilon^2} \rceil$ largest such value in $M$. This is sufficient to show that with high probability, the value of $\rho$ returned by Algorithm~\ref{alg:audit} is not an underestimate, in the sense that the set of points $S$ returned by the algorithm supports a value of $\rho$ at least so large.    
\end{proof}

% END COMMENT 
\fi







